\par\Needspace{6\baselineskip}\medskip
\SourceLeaf{157}\pdfbookmark[2]{112.}{article-112}
\begin{center}112.\end{center}
Let $x''$, $y''$, and $z''$ be the coordinates of a third position, $u''$ its argument of latitude, and $r''$ its radius vector. We denote the quantities
\[
r'r''\sin(u''-u'),\qquad rr''\sin(u''-u),\qquad rr'\sin(u'-u),
\]
which are twice the areas of the triangles enclosed by the second and third radius vectors, the first and third, and the first and second, by $n$, $n'$, and $n''$, respectively. We then have expressions for $x''$, $y''$, and $z''$ like those given in the preceding article for $x$, $y$, and $z$ and for $x'$, $y'$, and $z'$. With the aid of Lemma I of Article 78, the following equations are readily derived.
\[
\begin{aligned}
0&=nx-n'x'+n''x''\\
0&=ny-n'y'+n''y''\\
0&=nz-n'z'+n''z''.
\end{aligned}
\]
Let $\alpha$, $\alpha'$, and $\alpha''$ be the geocentric longitudes of the celestial body corresponding to those three positions, $\beta$, $\beta'$, and $\beta''$ its geocentric latitudes, and $\delta$, $\delta'$, and $\delta''$ its distances from the Earth projected onto the ecliptic. Further, let $L$, $L'$, and $L''$ be the corresponding heliocentric longitudes of the Earth and $B$, $B'$, and $B''$ its latitudes. We do not set these latitudes $=0$, so that parallax can be taken into account and any other plane can be adopted in place of the ecliptic, if desired. Finally, let $D$, $D'$, and $D''$ be the Earth's distances from the Sun projected onto the ecliptic. If $x$, $y$, and $z$ are now expressed in terms of $L$, $B$, $D$, $\alpha$, $\beta$, and $\delta$, and the coordinates pertaining to the second and third positions are expressed in the same way, the preceding equations take the following form.
\[
\begin{aligned}
1)\quad 0&=n(\delta\cos\alpha+D\cos L)-n'(\delta'\cos\alpha'+D'\cos L')\\
&\qquad+n''(\delta''\cos\alpha''+D''\cos L'')\\[.5ex]
2)\quad 0&=n(\delta\sin\alpha+D\sin L)-n'(\delta'\sin\alpha'+D'\sin L')\\
&\qquad+n''(\delta''\sin\alpha''+D''\sin L'')\\[.5ex]
3)\quad 0&=n(\delta\operatorname{tang}\beta+D\operatorname{tang}B)-n'(\delta'\operatorname{tang}\beta'+D'\operatorname{tang}B')\\
&\qquad+n''(\delta''\operatorname{tang}\beta''+D''\operatorname{tang}B'').
\end{aligned}
\]
If $\alpha$, $\beta$, $D$, $L$, and $B$, together with the analogous quantities for the other two positions, are regarded as known, and the equations are divided by $n$, $n'$, or $n''$, five unknowns remain. Two of these can therefore be eliminated, or the remaining three can be determined from any two of these unknowns. In this way the three equations lead to many important conclusions, of which we shall develop a few particularly noteworthy ones here.
